% ID: MAT-GAN-ELL-003
% Titolo: Semiellisse traslata: dominio e punti notevoli
% Disciplina: Matematica
% Area: geometria_analitica
% Argomento: ellisse
% Classe: Terza o quarta secondaria di secondo grado
% Difficolta: 4
% Tipo: grafico
% Tag: matematica, geometria_analitica, ellisse
% Fonte: verifica_ellisse_anonimizzata
% Provenienza: materiale fornito dal docente; testo anonimizzato; soluzione redatta per Exergo

\begin{esercizio}
Traccia il grafico della seguente funzione:
    $f(x)=3\sqrt{-x^2+4x}$
    avendo cura di segnalare le coordinate di tutti i punti notevoli che identificano la conica.
\end{esercizio}

\begin{soluzione}
Completando il quadrato, $y=3\sqrt{4-(x-2)^2}$, con $0\le x\le4$ e $y\ge0$.
Elevando al quadrato: $(x-2)^2/4+y^2/36=1$, ma si conserva solo la metà superiore.
Centro della conica completa: $(2,0)$; semiassi 2 orizzontale e 6 verticale;
estremi dell'arco $(0,0),(4,0)$ e massimo $(2,6)$.
I fuochi della conica completa sono $(2,\pm4\sqrt2)$: non sono punti dell'arco.
\end{soluzione}
